% SOURCE: original Chinese PDF, viewer pages 4--6.
\section{Linear Vector Spaces}
From the preceding section, the projection $s_n$ of the electron's spin angular momentum along any direction $n$ can take only two values. We may regard $s_n$ as a two-dimensional vector. To establish the mathematical description of quantum mechanics, we first discuss linear vector spaces.
\subsection*{1) Three-dimensional vector space}
\begin{center}\begin{tikzpicture}[>=Latex,scale=.85]
\coordinate (O) at (0,0); \draw[->] (O)--(2.0,.1)node[right]{$\vect e_2$};\draw[->](O)--(-1.25,-.8)node[left]{$\vect e_1$};\draw[->](O)--(0,1.8)node[above]{$\vect e_3$};\draw[->,very thick](O)--(1.05,1.6)node[right]{$\vect a$};
\end{tikzpicture}\end{center}
Arbitrary vector: $\vect a$.

Basis vectors: $\vect e_n$, $n=1,2,3$.

Completeness of the basis:
\begin{equation}
\vect a=\sum_{n=1}^{3}a_n\vect e_n.
\end{equation}
Inner product:
\begin{equation}
\vect a\cdot\vect b=\sum_{n,m}a_nb_m\vect e_n\cdot\vect e_m.
\end{equation}
Squared norm of a vector: $\vect a\cdot\vect a\geq0$.

If the basis vectors are orthonormal, $\vect e_n\cdot\vect e_m=\delta_{nm}$, the inner product has the matrix form
\begin{equation}
\vect a\cdot\vect b=\sum_n a_nb_n=\widetilde a\,b,
\quad b=\begin{pmatrix}b_1\\b_2\\b_3\end{pmatrix},\quad
 a=\begin{pmatrix}a_1\\a_2\\a_3\end{pmatrix},\quad
 \widetilde a=(a_1,a_2,a_3),
\end{equation}
where $\widetilde a$ is the transpose of $a$.

Components (matrix elements) of a vector:
\begin{equation}
a_n=\vect e_n\cdot\vect a.
\end{equation}
The vector $\vect a$ is its abstract or general form. The matrix $a$ is the representation of $\vect a$ in a specific coordinate system (representation). Its matrix elements depend on the choice of basis vectors; for example, Cartesian and spherical coordinate systems give different representations.

Operations on a vector, such as translations and rotations (operators), give another vector in the same three-dimensional space: $\op T\vect a=\vect b$.
\subsection*{2) Hilbert space}
Extend the three-dimensional vector space to a complex vector space of arbitrary dimension:
\begin{equation}
\text{three dimensions}\longrightarrow\text{arbitrary finite, infinite, or continuous dimension},
\end{equation}
\begin{equation}
\text{constant vectors}\longrightarrow\text{complex-function vectors}.
\end{equation}
Use Dirac notation (a ket) to denote an arbitrary vector: $\ket a$.

For a complex vector, introduce the bra $\bra a$ to represent its complex-conjugate vector. The bra and ket are not independent; they are conjugate to each other:
\begin{equation}
\ket a\longleftrightarrow\bra a.
\end{equation}
A vector can be represented by either a ket $\ket a$ or a bra $\bra a$.

In a complex-function vector space, constants are generally complex, so
\begin{equation}
\alpha\ket a\longleftrightarrow\bra a\alpha^*.
\end{equation}
For an operation (operator) $\op T$ in the complex space,
\begin{equation}
\op T\ket a\longleftrightarrow\bra a\op T^\dagger.
\end{equation}
$\op T^\dagger$ is the left operator corresponding to the right operator $\op T$, called the Hermitian-conjugate operator. Note that
\begin{equation}
\op T\op F\ket a=\op T(\op F\ket a)
\longleftrightarrow \bra a\op F^\dagger\op T^\dagger.
\end{equation}
We now choose a particular representation, taking an $N$-dimensional discrete space as an example.

Basis vectors: $\ket n$, $n=1,2,\ldots$.

Expansion in a complete basis (concrete representation of a vector):
\begin{equation}
\ket a=\sum_n a_n\ket n,\qquad \bra a=\sum_n a_n^*\bra n.
\end{equation}
Inner product:
\begin{equation}
\braket a b=\sum_{n,m}a_n^*b_m\braket n m,
\end{equation}
which is a complex number.

Squared norm: $\braket a a\geq0$. (Only when $\ket a$ and $\bra a$ are defined as conjugates can the squared norm be guaranteed nonnegative.)

Normalized vector: defining $\ket{\widetilde a}=\ket a/\sqrt{\braket a a}$ gives $\braket{\widetilde a}{\widetilde a}=1$; the vector is said to be normalized.

Orthonormality of basis vectors: $\braket n m=\delta_{nm}$.

Matrix form of the inner product:
\begin{equation}
\braket a b=\sum_n a_n^*b_n=a^\dagger b,
\qquad b=\begin{pmatrix}b_1\\b_2\\\vdots\\b_N\end{pmatrix},\quad
 a=\begin{pmatrix}a_1\\a_2\\\vdots\\a_N\end{pmatrix},
\end{equation}
where the Hermitian conjugate (complex-conjugate transpose) of the column matrix $a$ is the row matrix
\begin{equation}
a^\dagger=(a_1^*,a_2^*,\ldots,a_N^*).
\end{equation}
Evidently,
\begin{equation}
\braket a b=\braket b a^*,\qquad
\mel a{\op T}b=\mel b{\op T^\dagger}a^*.
\end{equation}
Components (matrix elements) of a vector: multiplying $\ket a=\sum_n a_n\ket n$ from the left by $\bra m$ gives
\begin{equation}
a_m=\braket m a.
\end{equation}
Completeness of the basis:
\begin{equation}
\ket a=\sum_n a_n\ket n=\sum_n\ket n\braket n a
=\left(\sum_n\ket n\bra n\right)\ket a.
\end{equation}
Because $\ket a$ is arbitrary,
\begin{equation}
\boxed{\sum_n\ket n\bra n=\mathbf 1},
\end{equation}
which is the completeness relation for the basis.

\paragraph{Representation of an operator (matrix form).}
Let $\ket b=\op F\ket a$. Then
\begin{equation}
\braket m b=\mel m{\op F}a
=\sum_n\mel m{\op F}n\braket n a,
\qquad F_{mn}=\mel m{\op F}n,
\end{equation}
and therefore
\begin{equation}
b_m=\sum_n F_{mn}a_n.
\end{equation}
In matrix form, $b=Fa$. The matrix $F$ represents the operator $\op F$; it is a square matrix with elements $F_{mn}$.
\paragraph{Representation of the Hermitian-conjugate operator (left operator).}
From $\bra b=\bra a\op F^\dagger$,
\begin{equation}
\braket b m=\mel a{\op F^\dagger}m
=\sum_n\braket a n\mel n{\op F^\dagger}m,
\end{equation}
so
\begin{equation}
b_m^*=\sum_n (F^\dagger)_{nm}a_n^*,\qquad
b_m=\sum_n[(F^\dagger)_{nm}]^*a_n.
\end{equation}
Comparison yields
\begin{equation}
[(F^\dagger)_{nm}]^*=F_{mn},\qquad
(F^\dagger)_{nm}=F_{mn}^*,\qquad
F^\dagger=(\widetilde F)^*.
\end{equation}
Thus $F^\dagger$ is the Hermitian-conjugate matrix of $F$.
\paragraph{Outer product.}
The outer product is $\ket a\bra b$.

Its representation is a square matrix ($a$ is a column matrix and $b^\dagger$ a row matrix), so the outer product is an operator. In fact,
\begin{equation}
(\ket a\bra b)\ket c=\ket a(\braket b c),
\end{equation}
showing that $\ket a\bra b$ maps $\ket c$ to another vector parallel to $\ket a$. Hence $\ket a\bra b$ is indeed an operator. Its representation is a square matrix with elements
\begin{equation}
(\ket a\bra b)_{mn}=\mel m{\ket a\bra b}n
=\braket m a\braket b n=a_mb_n^*.
\end{equation}